\begin{definition}[Metric $1$-current $\mathcal{M}_1(E)$]
  Let $E$ be a metric space equipped with a compatible Borel $\sigma$-algebra. A \emph{metric
  $1$-current} on $E$ is a structure consisting of a functional $T \colon \mathcal{D}^1(E) \to
  \mathbb{R}$ (\noderef{Form1}) together with:
  \begin{itemize}
    \item \emph{Additivity and homogeneity in $f$.} For all $\omega_1,\omega_2,\omega_{12} \in
      \mathcal{D}^1(E)$ with $\omega_1.\pi = \omega_2.\pi = \omega_{12}.\pi$ and $\omega_{12}.f =
      \omega_1.f + \omega_2.f$ (pointwise), $T(\omega_{12}) = T(\omega_1) + T(\omega_2)$; and for
      all $\omega,\omega_c \in \mathcal{D}^1(E)$ with $\omega_c.\pi = \omega.\pi$ and $c \in
      \mathbb{R}$ with $\omega_c.f = c\cdot\omega.f$ (pointwise), $T(\omega_c) = c\cdot T(\omega)$.
    \item \emph{Additivity and homogeneity in $\pi$.} The same two conditions with the roles of
      $f$ and $\pi$ exchanged: additivity when $\omega_1.f=\omega_2.f=\omega_{12}.f$ and
      $\omega_{12}.\pi = \omega_1.\pi + \omega_2.\pi$ pointwise, and homogeneity when
      $\omega_c.f=\omega.f$ and $\omega_c.\pi = c\cdot\omega.\pi$ pointwise.
    \item \emph{Continuity with respect to pointwise convergence of $\pi$.} For fixed $f$ (with
      any bound and Lipschitz witnesses) and $\pi$ (Lipschitz constant $\mathrm{Lip}(\pi)$), and
      any sequence $(\pi_n)$ with Lipschitz constants $(\mathrm{Lip}(\pi_n))$ uniformly bounded by
      some $L_0$ and $\pi_n(x) \to \pi(x)$ for every $x \in E$,
      \[
        T(f,\pi_n) \to T(f,\pi) \quad \text{as } n \to \infty
        .
      \]
    \item \emph{Locality.} For every $\omega=(f,\pi) \in \mathcal{D}^1(E)$, if there is an open set
      $U \supseteq \{x \in E : f(x) \neq 0\}$ on which $\pi$ is constant, then $T(\omega) = 0$.
    \item \emph{Finite mass.} $\mathbb{M}(T) := \mathrm{massOfFunctional}(T) \neq \infty$
      (\noderef{massOfFunctional}).
  \end{itemize}
  This is Ambrosio--Kirchheim's definition of a $k$-dimensional metric current
  \cite{AmbrosioKirchheim}*{Definition~3.1}, specialized to $k=1$ (paper.tex lines 1096--1109, which introduce
  $\mathcal{D}^1(E)$ and $\mathcal{M}_1(E)$ and cite \cite{AmbrosioKirchheim} for the precise
  definition rather than spelling it out). We record the transcription here as the external
  definitional input it is.

  \paragraph{Scope.} This node builds only the $k=1$ case (single $\pi$-argument), matching the
  needs of paper Theorem~4.4 (target \texttt{bbassertion}); it is not a general $k$-current.

  \paragraph{Why continuity with respect to $f$ is \emph{not} an axiom here.} Paper.tex line 1393
  states this explicitly: ``we recall that by definition, a metric $1$-current $T$ is continuous
  with respect to the pointwise convergence [of $\pi$], while the continuity with respect to [the
  analogous convergence of $f$] follows from [the mass-measure inequality] and Lebesgue's
  dominated convergence theorem.'' That is, continuity in $f$ is a \emph{consequence} of
  \emph{finite mass} together with additivity/homogeneity in $f$ (to reduce $T(f_n,\pi)-T(f,\pi)$
  to $T(f_n-f,\pi)$) and dominated convergence (to bound $\int_E |f_n-f|\,d\mu$ by the mass measure
  once $|f_n-f|\le 2\sup_n\|f_n\|_\infty$ and $f_n-f\to 0$ pointwise). Axiomatizing it directly, in
  addition to the finite-mass condition, would strengthen the definition of $\mathcal{M}_1(E)$
  beyond what the paper justifies, so it is omitted as a field and left to be derived at whichever
  node needs it.

  \paragraph{Why locality is kept even though it is unused elsewhere in the paper.} A search of the
  source shows no other use of the word ``local'' in this paper's argument. We nonetheless include
  it, because it is part of the actual mathematical content of ``$T \in \mathcal{M}_1(E)$'' in the
  Ambrosio--Kirchheim sense that paper.tex line 1101 invokes, and including it only
  \emph{strengthens} the hypothesis available to later nodes (a stronger, paper-faithful
  hypothesis on $T$ is safe to carry even when unused), rather than weakening any conclusion drawn
  from it.

  \paragraph{Modeling choice: bilinearity via matching side-conditions (interface decision, not a
  deviation).} \noderef{Form1} bundles a specific witnessing bound $\|f\|_\infty$ and specific
  witnessing Lipschitz constants $\mathrm{Lip}(f),\mathrm{Lip}(\pi)$ as data alongside the
  functions $f,\pi$ themselves, rather than only asserting their existence. Consequently the
  additivity/homogeneity conditions above are stated by universally quantifying over
  \emph{any} triple (or pair) of $\mathcal{D}^1(E)$-elements whose function-components stand in
  the stated pointwise relationship, rather than by explicitly constructing a ``sum'' or
  ``scalar multiple'' $1$-form as a new piece of data. This changes no mathematical content: since
  every valid combination of functions with matching pointwise relationships extends to \emph{some}
  witnessed $\mathcal{D}^1(E)$-element (any bound/Lipschitz overestimate is a valid witness, as
  \noderef{Form1}'s own modeling note observes), the universally-quantified form captures exactly
  the classical bilinearity condition on the underlying functions, for every witnessed instance
  that could occur.
\end{definition}
